\documentclass[a4paper, 10pt]{article}
\usepackage{../CEDT-Assignment-style}

\begin{document}
\subject
\asgntitle{1}{}{Week 1}{6733172621 Patthadon Phengpinij}{Claude (for\,\LaTeX\,styling and grammar checking)}

% ================================================================================ %
%                                    Problem 01                                    %
% ================================================================================ %
\begin{problem}
\textbf{Quantum Computing Account}
\par Please go to IBM Quantum Experience and register an account, \url{https://quantum.cloud.ibm.com/}.
\end{problem}

\begin{solution}
	Logged in to IBM Quantum Experience and registered an account successfully.
	\fig*[0.9\linewidth]{images/p1-answer.png}{}{}
\end{solution}
% ================================================================================ %

% ================================================================================ %
%                                    Problem 02                                    %
% ================================================================================ %
\begin{problem}
\textbf{My First Quantum Circuit}
\par Try to create a circuit. Click on IBM Quantum Composer and drag the gates to form the following circuit.

\fig*[0.3\linewidth]{images/p2-question.png}{}{}

The first icon is a quantum gate called Hadamard gate and the second icon is a measurement.
The input is from the left ($q_0$).
Do not worry if you do not understand what it is at this moment.
But good job! You have created your first quantum circuit!
\end{problem}

\begin{solution}
	Finished creating the quantum circuit as instructed.
	\fig*[0.3\linewidth]{images/p2-answer.png}{}{}
\end{solution}
% ================================================================================ %

% ================================================================================ %
%                                    Problem 03                                    %
% ================================================================================ %
\begin{problem}
\textbf{My First Quantum Computing}
\par Submit the job. Click on the result when it is finished. You see some like this.
You may submit to a simulator or real quantum hardware.

\fig*[0.3\linewidth]{images/p3-question.png}{}{}

We started with a state in \( \ket{0} \) (a basis state).
After the \textit{Hadamard gate}, it becomes a superposition of \( \ket{0} \) and \( \ket{1} \), i.e. \( \sim \ket{0} + \ket{1} \) (we will learn later).
Then we performed a measurement on the output and the superposition state collapses to either \( \ket{0} \) or \( \ket{1} \).
The histogram shows that about 50\% of the chance we obtain \( \ket{0} \) and another 50\% of the chance we obtain \( \ket{1} \).
Cool! We have done our first quantum computation.
This job was submitted to a real quantum computer, called ``ibmq\_casablanca.''
It is a 5-qubit quantum computer. However, this computer is gone.
But definitely, you will find other more powerful ones.
\end{problem}

\begin{solution}
	The job was done successfully and the result was obtained as instructed.
	\fig*[0.7\linewidth]{images/p3-answer.png}{}{}
\end{solution}
% ================================================================================ %

% ================================================================================ %
%                                    Problem 04                                    %
% ================================================================================ %
\begin{problem}
\textbf{Google Colab and Python}
\par We need to do a lot of matrix operations. We will choose Google Colab, in which we can run Python easily.
Since IBM Quantum Experience also uses Python, it will save your efforts.
Please open a Google Colab account here \url{https://colab.research.google.com/}.
\end{problem}

\begin{solution}
	I decided to use my local computer to run the Jupyter notebook instead of Google Colab.
\end{solution}
% ================================================================================ %

% ================================================================================ %
%                                    Problem 05                                    %
% ================================================================================ %
\begin{problem}
\textbf{Vector Inner Product Using Python}
\par We use the NumPy library to perform vector and matrix operations. To import the library, type

\begin{codingbox}
import numpy as np
\end{codingbox}

Then let's implement
\[
	\vec{v_1} \cdot \vec{v_2}
	= \begin{pmatrix}\sqrt{3} & 1\end{pmatrix} \begin{pmatrix}1 \\ 1\end{pmatrix}
	= \paren{\sqrt{3} \times 1} + \paren{1 \times 1}
	= \sqrt{3} + 1
	\approx 2.732
\]

\begin{codingbox}
v1 = np.array([np.sqrt(3), 1])
v2 = np.array([1, 1])
inner_product = np.vdot(v1, v2)

print(inner_product)
\end{codingbox}

Type Shift-Enter to let it run. Do you see the same result?
Note that we use \texttt{np.vdot} instead of \texttt{np.dot} because we need to use complex vector dot products in quantum computing.
Try also \texttt{print(v1.shape)}, which will give you (2, 1).
It means we constructed a \( 2 \times 1 \) matrix, which is a column vector with 2 elements.
\end{problem}

\begin{solution}
	Try running the code in a Python environment and obtained the same result of approximately 2.732.

	\begin{codingbox}
import numpy as np

v1 = np.array([[np.sqrt(3)], [1]])
v2 = np.array([[1], [1]])
inner_product = np.vdot(v1, v2)

print(f"Inner product: {inner_product:.3f}")
print("v1 Shape:", v1.shape)
print("v2 Shape:", v2.shape)
\end{codingbox}

\end{solution}
% ================================================================================ %

% ================================================================================ %
%                                    Problem 06                                    %
% ================================================================================ %
\begin{problem}
\textbf{Inner Product with Complex Coefficient 1}
\par Now, let's try \( \vec{v_1} = \begin{pmatrix} i \\ 1 \end{pmatrix} \) and \( \vec{v_2} = \begin{pmatrix} -i \\ 1 \end{pmatrix} \).
Firstly, try to use hand calculation. Then use Python to check if it is correct.
Note that in the NumPy library, you need to put \( i \) as \( j \).
For example, \( 5 + i \) should be entered as \( 5 + 1j \) (putting the ``1'' is necessary).

Think about that before reading the code below:

\begin{codingbox}
v1 = np.array([[1j], [1]])
v2 = np.array([[-1j], [1]])
inner_result = np.vdot(v1, v2)
print(inner_result)
\end{codingbox}

\end{problem}

\begin{solution}
	First, try to calculate the inner product manually:
	\begin{align*}
		\vec{v_1} \cdot \vec{v_2} & = \begin{pmatrix} i \\ 1 \end{pmatrix} \cdot \begin{pmatrix} -i \\ 1 \end{pmatrix} \\
		                          & = \begin{pmatrix} i^* & 1^* \end{pmatrix} \begin{pmatrix} -i \\ 1 \end{pmatrix}    \\
		                          & = \begin{pmatrix} -i & 1 \end{pmatrix} \begin{pmatrix} -i \\ 1 \end{pmatrix}       \\
		                          & = (-i \times -i) + (1 \times 1)                                                    \\
		                          & = -1 + 1                                                                           \\
		\vec{v_1} \cdot \vec{v_2} & = 0
	\end{align*}

	Then, re-check the result using Python code given by the problem.
	We can see that the result is also 0, which is the same as what we calculated manually.
\end{solution}
% ================================================================================ %

% ================================================================================ %
%                                    Problem 07                                    %
% ================================================================================ %
\begin{problem}
\textbf{Inner Product with Complex Coefficient 2}
\par Now, let's try \( \vec{v_1} = \begin{pmatrix} 1 + 2i \\ 1 \end{pmatrix} \) and \( \vec{v_2} = \begin{pmatrix} 1 - i \\ 1 \end{pmatrix} \).
\end{problem}

\begin{solution}
	First, we calculate the inner product manually:
	\begin{align*}
		\vec{v_1} \cdot \vec{v_2} & = \begin{pmatrix} 1 + 2i \\ 1 \end{pmatrix} \cdot \begin{pmatrix} 1 - i \\ 1 \end{pmatrix} \\
		                          & = \begin{pmatrix} (1 + 2i)^* & 1^* \end{pmatrix} \begin{pmatrix} 1 - i \\ 1 \end{pmatrix}  \\
		                          & = \begin{pmatrix} 1 - 2i & 1 \end{pmatrix} \begin{pmatrix} 1 - i \\ 1 \end{pmatrix}        \\
		                          & = (1 - 2i)(1 - i) + (1)(1)                                                                 \\
		                          & = 1 - i - 2i + 2i^2 + 1                                                                    \\
		                          & = 1 - 3i - 2 + 1                                                                           \\
		\vec{v_1} \cdot \vec{v_2} & = -3i
	\end{align*}

	Then, re-check the result using the following Python code:

	\begin{codingbox}
v1 = np.array([[1 + 2j], [1]])
v2 = np.array([[1 - 1j], [1]])
inner_result = np.vdot(v1, v2)

print(f"Inner product: {inner_result:.3f}")
\end{codingbox}

	We can see that the result is also -3i, which is the same as what we calculated manually.
\end{solution}
% ================================================================================ %

% ================================================================================ %
%                                     Appendix                                     %
% ================================================================================ %

\myappendix{code/assignment-01-code.pdf}{Assignment 01 | Code Solution}

% ================================================================================ %

\end{document}